\documentclass[a4paper,12pt]{article}
\usepackage[utf8]{inputenc}
\usepackage[T1]{fontenc}
\usepackage[ngerman]{babel}
\usepackage{amsmath, amssymb}
\usepackage{geometry}
\usepackage{tikz}
\usetikzlibrary{shapes.geometric, arrows.meta, positioning, fit, backgrounds}
\usepackage{parskip}
\usepackage{titlesec}
\usepackage{booktabs}
\usepackage{enumitem}

\geometry{left=3cm, right=3cm, top=2.5cm, bottom=2.5cm}

\titleformat{\section}{\large\bfseries}{}{0em}{}[\titlerule]
\titleformat{\subsection}{\normalsize\bfseries}{}{0em}{}

\tikzset{
  box/.style={rectangle, draw, rounded corners=3pt, text width=4.5cm,
              align=center, minimum height=1.0cm, font=\small, inner sep=4pt},
  boxw/.style={box, text width=5.5cm},
  arrow/.style={-{Latex[length=2mm]}, thick},
}

\begin{document}

\begin{center}
  {\LARGE\bfseries Studierhinweise zu Lektion 4}\\[0.5em]
  {\large Mathematische Grundlagen (Modul 61111)}
\end{center}

\vspace{1em}

Diese Lektion eröffnet den Analysis-Teil des Moduls. Im Zentrum steht die Frage,
was die reellen Zahlen~$\mathbb{R}$ gegenüber anderen Körpern so besonders macht -- und
warum Analysis ohne $\mathbb{R}$ nicht funktioniert.

In \textbf{Kapitel~12} werden neun Axiome für $\mathbb{R}$ vorgestellt: fünf Körperaxiome,
drei Ordnungsaxiome und das entscheidende Schnittaxiom (Axiom der Ordnungsvollständigkeit).
Aus diesen Axiomen leiten wir systematisch die Bruchrechnung, den Umgang mit Ungleichungen,
den Absolutbetrag, Intervalle sowie das Supremumsprinzip her -- das Herzstück, das
$\mathbb{R}$ von~$\mathbb{Q}$ unterscheidet. Am Ende von Kapitel~12 zeigen wir, dass
$p$-te Wurzeln aus nichtnegativen reellen Zahlen stets existieren und eindeutig sind.

In \textbf{Kapitel~13} wird der Begriff des Grenzwertes einer reellen Folge
präzisiert. Dieser Begriff ist das unverzichtbare Werkzeug aller Analysis. Neben
der formalen $\varepsilon$-Definition lernen Sie Rechenregeln für konvergente Folgen
sowie vier wichtige Konvergenzprinzipien kennen: das Monotonieprinzip, den Satz
von Bolzano-Weierstraß, das Cauchy-Kriterium und das Prinzip der Intervallschachtelung.

Lesen Sie diesen Lehrtext \emph{aktiv}: Schreiben Sie jeden Beweisschritt selbst nach,
und überprüfen Sie jede Ungleichung an einem kleinen Zahlenbeispiel, bevor Sie weiterlesen.
Kapitel~13 ist Grundlage der gesamten weiteren Analysis -- lassen Sie hier keine Lücken.

% -----------------------------------------------------------------------
\section*{Struktur der Lektion 4}

\begin{center}
\begin{tikzpicture}[node distance=0.55cm and 0.9cm]

  % Column headers
  \node[font=\bfseries\small] (h1) at (0,0)    {Kap.\,12: Reelle Zahlen};
  \node[font=\bfseries\small] (h2) at (7.5,0)  {Kap.\,13: Grenzwerte von Folgen};

  % Kap 12 boxes
  \node[box, below=0.8cm of h1] (ax1) {12.1 Körperaxiome\\(A1)--(A5)};
  \node[box, below=0.5cm of ax1] (ax2) {12.1 Ordnungsaxiome\\(A6)--(A8)};
  \node[box, below=0.5cm of ax2] (ax3) {12.1 Schnittaxiom\\(A9)};
  \node[box, below=0.5cm of ax3] (fk)  {12.2.1 Bruchrechnung\\Vorzeichenregeln};
  \node[box, below=0.5cm of fk]  (un)  {12.2.2 Ungleichungen\\Betrag, Abstand, Intervalle};
  \node[box, below=0.5cm of un]  (sup) {12.2.3 Supremum, Infimum\\Supremumsprinzip};
  \node[box, below=0.5cm of sup] (wur) {12.2.4 $p$-te Wurzeln\\Binomischer Lehrsatz};

  % Kap 13 boxes
  \node[box, below=0.8cm of h2] (gw)  {13.1 Grenzwertbegriff\\$\varepsilon$-Definition, Limes};
  \node[box, below=0.5cm of gw]  (eig) {13.2 Eigenschaften\\Eindeutigkeit, Beschränktheit};
  \node[box, below=0.5cm of eig] (div) {13.3 Divergente Folgen};
  \node[box, below=0.5cm of div] (rec) {13.4 Rechenregeln\\Summe, Produkt, Quotient};
  \node[box, below=0.5cm of rec] (mon) {13.5 Monotonieprinzip};
  \node[box, below=0.5cm of mon] (bw)  {13.5 Bolzano-Weierstraß};
  \node[box, below=0.5cm of bw]  (cau) {13.5 Cauchy-Kriterium};
  \node[box, below=0.5cm of cau] (ivs) {13.5 Intervallschachtelung};

  % Arrows Kap 12
  \draw[arrow] (ax1) -- (ax2);
  \draw[arrow] (ax2) -- (ax3);
  \draw[arrow] (ax1.east) -- ++(0.3,0) |- (fk.east);
  \draw[arrow] (ax2.east) -- ++(0.3,0) |- (un.east);
  \draw[arrow] (ax3) -- (sup);
  \draw[arrow] (sup) -- (wur);

  % Arrows Kap 13
  \draw[arrow] (gw)  -- (eig);
  \draw[arrow] (eig) -- (div);
  \draw[arrow] (div) -- (rec);
  \draw[arrow] (rec) -- (mon);
  \draw[arrow] (mon) -- (bw);
  \draw[arrow] (bw)  -- (cau);
  \draw[arrow] (cau) -- (ivs);

  % Cross-chapter
  \draw[arrow] (sup.east) -- ++(0.4,0) |- (gw.west);
  \draw[arrow] (sup.east) -- ++(0.4,0) |- (mon.west);

\end{tikzpicture}
\end{center}

\newpage

% -----------------------------------------------------------------------
\section*{Zielelement 12.1 -- Die Axiome der reellen Zahlen}

\subsection*{Lerninhalte}

\begin{center}
\begin{tikzpicture}[node distance=0.55cm and 1.0cm]
  \node[box] (k)  {Körperaxiome (A1)--(A5):\\$\mathbb{R}$ ist ein Körper};
  \node[box, below=0.5cm of k] (o)  {Ordnungsaxiome (A6)--(A8):\\Trichotomie, Transitivität,\\Monotoniegesetze};
  \node[box, below=0.5cm of o] (s)  {Schnittaxiom (A9):\\Dedekind'scher Schnitt $(A|B)$,\\Trennungszahl};
  \node[box, below=0.5cm of s] (u)  {Unterschied $\mathbb{R}$ vs.\ $\mathbb{Q}$:\\angeordneter Körper -- nicht hinreichend};
  \draw[arrow] (k) -- (o);
  \draw[arrow] (o) -- (s);
  \draw[arrow] (s) -- (u);
\end{tikzpicture}
\end{center}

\subsection*{Lernziele}

Nach Durcharbeiten dieses Abschnitts sollten Sie:
\begin{itemize}
  \item die neun Axiome der reellen Zahlen benennen und in drei Gruppen einteilen können,
  \item erklären können, warum weder die Körperaxiome noch die Ordnungsaxiome allein $\mathbb{R}$
        von $\mathbb{Q}$ unterscheiden,
  \item den Begriff des Dedekind'schen Schnittes $(A|B)$ und der Trennungszahl definieren können,
  \item verstehen, warum $\mathbb{F}_2$ kein angeordneter Körper ist,
  \item das Schnittaxiom (A9) in eigenen Worten erläutern können (geometrische Intuition:
        Lückenlosigkeit der Zahlengeraden).
\end{itemize}

\subsection*{Selbstkontrollelement 12.1}

Sei $A = \{x \in \mathbb{Q} \mid x < \sqrt{2}\}$ und $B = \{x \in \mathbb{Q} \mid x > \sqrt{2}\}$.
Erklären Sie, warum $(A|B)$ ein Dedekind'scher Schnitt der rationalen Zahlen ist,
und was das Schnittaxiom für die reellen Zahlen in diesem Zusammenhang leistet.

\newpage

% -----------------------------------------------------------------------
\section*{Zielelement 12.2.1--12.2.2 -- Folgerungen: Bruchrechnung und Ungleichungen}

\subsection*{Lerninhalte}

\begin{center}
\begin{tikzpicture}[node distance=0.55cm and 1.0cm]
  \node[box] (br)  {Bruchrechnung, Vorzeichen\\ (Prop.\,12.2.1)};
  \node[box, below=0.5cm of br] (un) {Grundlegendes über Ungleichungen\\ (Prop.\,12.2.3--12.2.5)};
  \node[box, below=0.5cm of un] (be) {Bernoulli'sche Ungleichung $(1+x)^n \geq 1+nx$};
  \node[box, below=0.5cm of be] (bt) {Betrag $|a|$, Dreiecksungleichung\\$|a+b| \leq |a|+|b|$};
  \node[box, below=0.5cm of bt] (ab) {Abstand $d(a,b) = |a-b|$\\$\varepsilon$-Umgebungen $U_\varepsilon(x_0)$};
  \node[box, below=0.5cm of ab] (iv) {Intervalle: offen, abgeschlossen,\\ halboffen, uneigentlich};
  \draw[arrow] (br) -- (un);
  \draw[arrow] (un) -- (be);
  \draw[arrow] (be) -- (bt);
  \draw[arrow] (bt) -- (ab);
  \draw[arrow] (ab) -- (iv);
\end{tikzpicture}
\end{center}

\subsection*{Lernziele}

Nach Durcharbeiten dieses Abschnitts sollten Sie:
\begin{itemize}
  \item Ungleichungen mit den Grundregeln (Addition, Multiplikation, Vorzeichenwechsel)
        manipulieren und das Umdrehen des Ungleichungszeichens bei Multiplikation mit negativen
        Zahlen sicher anwenden können,
  \item die Bernoulli'sche Ungleichung anwenden und ihren Beweis per Induktion skizzieren können,
  \item den Absolutbetrag und seine Eigenschaften (Dreiecksungleichung, Multiplikativität) kennen,
  \item $\varepsilon$-Umgebungen als Intervallschreibweise und als Betragsbedingung
        $|x - x_0| < \varepsilon$ ausdrücken können,
  \item die verschiedenen Intervalltypen $(a,b)$, $[a,b]$, $(a,b]$, $[a,b)$ sowie
        uneigentliche Intervalle benennen und auf der Zahlengeraden darstellen können.
\end{itemize}

\subsection*{Selbstkontrollelement 12.2.1--12.2.2}

Beweisen Sie: Für alle $a, b \in \mathbb{R}$ gilt $|a - b| \geq \bigl||a| - |b|\bigr|$.\\
\emph{Hinweis:} Wenden Sie die Dreiecksungleichung auf $a = (a-b)+b$ an.

\newpage

% -----------------------------------------------------------------------
\section*{Zielelement 12.2.3 -- Supremum, Infimum und Supremumsprinzip}

\subsection*{Lerninhalte}

\begin{center}
\begin{tikzpicture}[node distance=0.55cm and 1.0cm]
  \node[box] (mm) {Minimum, Maximum\\(müssen in $M$ liegen)};
  \node[box, below=0.5cm of mm] (sc) {Obere und untere Schranken\\(müssen nicht in $M$ liegen)};
  \node[box, below=0.5cm of sc] (si) {Supremum $\sup M$\\Infimum $\inf M$};
  \node[box, below=0.5cm of si] (sp) {Supremumsprinzip:\\jede nicht leere, nach oben beschränkte\\Menge besitzt ein Supremum};
  \node[box, below=0.5cm of sp] (ae) {Satz des Archimedes:\\$\mathbb{N}$ ist nicht beschränkt};
  \node[box, below=0.5cm of ae] (eu) {Satz des Eudoxos:\\$\forall\varepsilon > 0\;\exists\,m\in\mathbb{N}: \tfrac{1}{m} < \varepsilon$};
  \node[box, below=0.5cm of eu] (dq) {$\mathbb{Q}$ liegt dicht in $\mathbb{R}$};
  \draw[arrow] (mm) -- (sc);
  \draw[arrow] (sc) -- (si);
  \draw[arrow] (si) -- (sp);
  \draw[arrow] (sp) -- (ae);
  \draw[arrow] (ae) -- (eu);
  \draw[arrow] (eu) -- (dq);
\end{tikzpicture}
\end{center}

\subsection*{Lernziele}

Nach Durcharbeiten dieses Abschnitts sollten Sie:
\begin{itemize}
  \item den Unterschied zwischen Minimum/Maximum und unteren/oberen Schranken erklären können,
  \item Supremum und Infimum einer Menge $M$ anhand der Definition 12.2.47 identifizieren können,
  \item das Supremumsprinzip (Satz 12.2.52) formulieren und seine Äquivalenz zum Schnittaxiom
        erläutern können (Korollar 12.2.56),
  \item den Satz des Archimedes (12.2.57) und den Satz des Eudoxos (12.2.58) anwenden können,
  \item verstehen, warum $\mathbb{Q}$ dicht in $\mathbb{R}$ liegt (Korollar 12.2.59).
\end{itemize}

\subsection*{Selbstkontrollelement 12.2.3}

Sei $M = \bigl\{\tfrac{n-1}{n+1} \bigm| n \in \mathbb{N}\bigr\}$.
Bestimmen Sie $\sup M$, $\inf M$, $\max M$ (falls existent) und $\min M$ (falls existent)
und begründen Sie Ihre Antwort.

\newpage

% -----------------------------------------------------------------------
\section*{Zielelement 12.2.4 -- $p$-te Wurzeln}

\subsection*{Lerninhalte}

\begin{center}
\begin{tikzpicture}[node distance=0.55cm and 1.0cm]
  \node[box] (bk)  {Binomialkoeffizient $\binom{n}{k}$};
  \node[box, below=0.5cm of bk]  (bl) {Binomischer Lehrsatz\\$(a+b)^p = \sum_{k=0}^p \binom{p}{k}a^{p-k}b^k$};
  \node[box, below=0.5cm of bl]  (de) {Definition $p$-te Wurzel aus $a\geq 0$:\\$r \geq 0$ mit $r^p = a$};
  \node[box, below=0.5cm of de]  (eu) {Eindeutigkeit (Lemma 12.2.66):\\$x < y \Leftrightarrow x^p < y^p$};
  \node[box, below=0.5cm of eu]  (ex) {Existenz (Satz 12.2.68):\\Supremumsprinzip liefert $\sqrt[p]{a}$};
  \node[box, below=0.5cm of ex]  (po) {Potenzen mit rationalem Exponent\\$a^{s/p} = \sqrt[p]{a^s}$, Potenzregeln};
  \draw[arrow] (bk) -- (bl);
  \draw[arrow] (bl) -- (de);
  \draw[arrow] (de) -- (eu);
  \draw[arrow] (eu) -- (ex);
  \draw[arrow] (ex) -- (po);
\end{tikzpicture}
\end{center}

\subsection*{Lernziele}

Nach Durcharbeiten dieses Abschnitts sollten Sie:
\begin{itemize}
  \item den Binomialkoeffizienten berechnen und den binomischen Lehrsatz anwenden können,
  \item die Definition der $p$-ten Wurzel kennen und verstehen, warum nur $r \geq 0$
        zugelassen ist,
  \item die Eindeutigkeit von $\sqrt[p]{a}$ mit Lemma~12.2.66 begründen können,
  \item die Beweisidee der Existenz (Satz 12.2.68) skizzieren können (Supremum
        der Menge $M = \{y \geq 0 \mid y^p < a\}$),
  \item die Potenzregeln für rationale Exponenten (Proposition 12.2.78) kennen,
  \item wissen, warum $\sqrt{2}$ irrational ist (Proposition 12.2.72).
\end{itemize}

\subsection*{Selbstkontrollelement 12.2.4}

Berechnen Sie $\binom{6}{2}$ und $\binom{6}{4}$.
Geben Sie außerdem an: Warum kann $\sqrt[3]{-8}$ gemäß Definition 12.2.64 nicht als dritte Wurzel aus $-8$ bezeichnet werden?

\newpage

% -----------------------------------------------------------------------
\section*{Zielelement 13.1--13.2 -- Der Grenzwertbegriff}

\subsection*{Lerninhalte}

\begin{center}
\begin{tikzpicture}[node distance=0.55cm and 1.0cm]
  \node[box] (fo) {Reelle Folge $(a_n)$:\\Abbildung $f:\mathbb{N}\to\mathbb{R}$};
  \node[box, below=0.5cm of fo] (fa) {„Fast alle" Folgenglieder,\\Endstücke $(a_n)_{n>n_0}$};
  \node[box, below=0.5cm of fa] (ko) {Konvergenz gegen $a$:\\$\forall\varepsilon>0\;\exists\,n_0:\forall n>n_0:\;|a_n - a|<\varepsilon$};
  \node[box, below=0.5cm of ko] (gw) {Grenzwert (Limes) $\lim_{n\to\infty} a_n = a$\\Eindeutigkeit (Prop.\,13.2.1)};
  \node[box, below=0.5cm of gw] (tf) {Teilfolge $(a_{n_k})$:\\Konvergenz gegen denselben Grenzwert};
  \node[box, below=0.5cm of tf] (bs) {Beschränktheit konvergenter Folgen};
  \draw[arrow] (fo) -- (fa);
  \draw[arrow] (fa) -- (ko);
  \draw[arrow] (ko) -- (gw);
  \draw[arrow] (gw) -- (tf);
  \draw[arrow] (tf) -- (bs);
\end{tikzpicture}
\end{center}

\subsection*{Lernziele}

Nach Durcharbeiten dieses Abschnitts sollten Sie:
\begin{itemize}
  \item die $\varepsilon$-Definition der Konvergenz (Definition 13.1.8) auswendig kennen
        und auf konkrete Folgen anwenden können,
  \item verstehen, was „fast alle Folgenglieder" bedeutet, und den Unterschied
        zu „alle" erläutern können,
  \item die Konvergenz einfacher Folgen (z.\,B.\ $(\tfrac{1}{n})$, $(\tfrac{1}{\sqrt[p]{n}})$,
        $(q^n)$ für $|q|<1$) direkt aus der Definition beweisen können,
  \item die Eindeutigkeit des Grenzwertes (Proposition 13.2.1) kennen,
  \item verstehen, warum jede Teilfolge einer konvergenten Folge gegen denselben
        Grenzwert konvergiert, und dieses Kriterium zur Divergenzprüfung einsetzen können.
\end{itemize}

\subsection*{Selbstkontrollelement 13.1--13.2}

Sei $(a_n) = \bigl(\tfrac{2n-14}{n+1}\bigr)$. Bestimmen Sie den Grenzwert und zeigen Sie
mit der $\varepsilon$-Definition, dass die Folge gegen diesen Grenzwert konvergiert.

\newpage

% -----------------------------------------------------------------------
\section*{Zielelement 13.3--13.4 -- Divergenz und Rechnen mit Folgen}

\subsection*{Lerninhalte}

\begin{center}
\begin{tikzpicture}[node distance=0.55cm and 1.0cm]
  \node[box] (dv)  {Divergenz: Negation\\der Konvergenz\\(Proposition 13.3.2)};
  \node[box, below=0.5cm of dv] (vg) {Vergleichssatz (Prop.\,13.4.1):\\$a_n \leq b_n \Rightarrow \lim a_n \leq \lim b_n$};
  \node[box, below=0.5cm of vg] (es) {Einschnürungssatz (Prop.\,13.4.3):\\$a_n \leq c_n \leq b_n \Rightarrow \lim c_n = \lim a_n$};
  \node[box, below=0.5cm of es] (nf) {Nullfolge $\times$ beschränkte Folge\\= Nullfolge (Prop.\,13.4.9)};
  \node[box, below=0.5cm of nf] (rr) {Rechenregeln (Prop.\,13.4.12):\\$\lim(a_n \pm b_n)$, $\lim a_n b_n$,\\$\lim(a_n / b_n)$};
  \draw[arrow] (dv) -- (vg);
  \draw[arrow] (vg) -- (es);
  \draw[arrow] (es) -- (nf);
  \draw[arrow] (nf) -- (rr);
\end{tikzpicture}
\end{center}

\subsection*{Lernziele}

Nach Durcharbeiten dieses Abschnitts sollten Sie:
\begin{itemize}
  \item Divergenz einer Folge formal charakterisieren und an Beispielen nachweisen können
        (z.\,B.\ $((-1)^n)$, $(n)$),
  \item Vergleichssatz und Einschnürungssatz anwenden können,
  \item die Rechenregeln für konvergente Folgen (Proposition 13.4.12) anwenden und dabei
        beachten, dass sie nur von rechts nach links gelesen werden dürfen,
  \item Grenzwerte rationaler Folgen durch Kürzen des führenden Terms berechnen können,
  \item wissen, dass aus der Konvergenz von $(a_n b_n)$ \emph{nicht} auf die Konvergenz
        von $(a_n)$ oder $(b_n)$ geschlossen werden darf.
\end{itemize}

\subsection*{Selbstkontrollelement 13.3--13.4}

Berechnen Sie $\displaystyle\lim_{n\to\infty}\frac{n^3+2n-5}{5n^3+n^2}$ und begründen Sie
jeden Schritt mit den Rechenregeln.
Geben Sie außerdem ein Beispiel für eine Nullfolge $(b_n)$ mit $b_n \neq 0$, für die
$\bigl(\tfrac{1}{b_n}\bigr)$ divergiert.

\newpage

% -----------------------------------------------------------------------
\section*{Zielelement 13.5 -- Vier Konvergenzprinzipien}

\subsection*{Lerninhalte}

\begin{center}
\begin{tikzpicture}[node distance=0.55cm and 1.0cm]
  \node[box] (mp) {Monotonieprinzip (Satz 13.5.2):\\monotone $+$ beschränkte Folge\\$\Rightarrow$ konvergent};
  \node[box, right=0.7cm of mp] (bw) {Bolzano-Weierstraß (Satz 13.5.13):\\beschränkte Folge $\Rightarrow$\\konvergente Teilfolge};
  \node[box, below=0.8cm of mp] (cf) {Cauchyfolge (Def.\,13.5.14):\\$\forall\varepsilon>0\;\exists\,n_0:\forall m,n>n_0:\;|a_m-a_n|<\varepsilon$};
  \node[box, right=0.7cm of cf] (cp) {Cauchy-Prinzip (Satz 13.5.16):\\konvergent $\Leftrightarrow$ Cauchyfolge};
  \node[box, below=0.8cm of cf] (is) {Intervallschachtelung $\langle a_n|b_n\rangle$\\(Satz 13.5.22):\\genau eine Zahl in allen $[a_n,b_n]$};
  \node[box, right=0.7cm of is] (eu) {Euler'sche Zahl $e$:\\$e = \lim_{n\to\infty}(1+\tfrac{1}{n})^n$};
  \draw[arrow] (mp.south) -- (cf.north);
  \draw[arrow] (bw.south) -- (cp.north);
  \draw[arrow] (cf) -- (is);
\end{tikzpicture}
\end{center}

\subsection*{Lernziele}

Nach Durcharbeiten dieses Abschnitts sollten Sie:
\begin{itemize}
  \item das Monotonieprinzip formulieren, seinen Beweis skizzieren und anwenden können
        (insbesondere: monotone Folge konvergiert $\Leftrightarrow$ beschränkt),
  \item den Satz von Bolzano-Weierstraß (jede beschränkte Folge hat eine konvergente Teilfolge)
        kennen und zum Nachweis von Konvergenzeigenschaften verwenden können,
  \item den Begriff der Cauchyfolge erklären und Cauchy-Folgen erkennen,
  \item das Cauchy-Konvergenzprinzip (konvergent $\Leftrightarrow$ Cauchyfolge) anwenden können,
  \item das Prinzip der Intervallschachtelung verstehen und wissen, dass genau eine reelle Zahl
        in allen Intervallen liegt,
  \item die Euler'sche Zahl $e$ als Grenzwert von $\bigl((1+\tfrac{1}{n})^{n+1}\bigr)$
        und $\bigl(\sum_{k=0}^{n}\tfrac{1}{k!}\bigr)$ kennen.
\end{itemize}

\subsection*{Selbstkontrollelement 13.5}

Sei $(a_n)$ mit $a_n = \sum_{k=1}^{n}\tfrac{1}{k^2}$. Zeigen Sie, dass $(a_n)$ monoton wachsend
und beschränkt ist (Hinweis: $\tfrac{1}{k^2} \leq \tfrac{1}{k(k-1)}$ für $k \geq 2$), und
schließen Sie auf Konvergenz. Ist $\lim a_n$ explizit bekannt?

\end{document}
